%% ma: L10-B14-0030
%% lop: 10
%% chuong: 5
%% bai: 14
%% dang: 10.14.5
%% loai: DS
%% muc: [NB, TH, TH, VD]
%% dap_an: "ĐSĐS"
%% nguon: "chuyen-de-cac-so-dac-trung-cua-mau-so-lieu-khong-ghep-nhom-toan-10-chuong-trinh-moi.pdf (GV. Phan Nhật Linh) (P04), Bài 03 Dạng 1 Phần II Câu 7"
%% kien_thuc: "Lớp 10 Bài 13 (số trung bình, tứ phân vị), Bài 14 (khoảng biến thiên, khoảng tứ phân vị, giá trị bất thường)"
%% trang_thai: chinh-thuc
%% ghi_chu: "Tự duyệt theo yêu cầu giáo viên 22/10/2026."
\begin{ex}
Số điểm mà các vận động viên bóng rổ ghi được trong một trận đấu như sau:
\begin{center}
\begin{tabular}{|c|c|c|c|c|c|c|c|}\hline
$9$ & $8$ & $15$ & $14$ & $11$ & $20$ & $13$ & $17$\\\hline
\end{tabular}
\end{center}
Xét tính đúng sai của các khẳng định sau:
\choiceTF{\True Khoảng biến thiên của mẫu số liệu là $12$}{Số trung bình của mẫu số liệu là $13{,}5$}{\True Khoảng tứ phân vị của mẫu số liệu là $6$}{Mẫu số liệu có giá trị bất thường}
\loigiai{Sắp xếp: $8$; $9$; $11$; $13$; $14$; $15$; $17$; $20$.

a) Đúng: $R=20-8=12$.

b) Sai: $\overline x=\dfrac{107}8=13{,}375$.

c) Đúng: $Q_1=\dfrac{9+11}2=10$; $Q_3=\dfrac{15+17}2=16$; $\Delta_Q=6$.

d) Sai: $Q_1-1{,}5\Delta_Q=1$; $Q_3+1{,}5\Delta_Q=25$; không có giá trị nào nằm ngoài đoạn $[1;25]$.}
\end{ex}
